\section{Ergodic Gibbs measures}

Throughout, $\Theta$ is a countable subgroup of the transformation group $T$ of $\Cfg$, $\mathcal I$
its invariant $\sigma$-algebra, $\Prob_\Theta$ the $\Theta$-invariant random fields and
$\mathcal G_\Theta(\gamma) = \mathcal G(\gamma) \cap \Prob_\Theta$.

\begin{theorem}[Georgii (14.10)]
    \label{thm:erg-1410}
    \uses{def:erg-invariant, def:erg-invariant-events, thm:erg-145, prop:app-pa-kernel}
    \lean{MeasureTheory.GibbsMeasure.exists_isPAKernel_invariantFields_shiftGroup, MeasureTheory.GibbsMeasure.exists_isPAKernel_invariantFields_shiftGroup_int, MeasureTheory.GibbsMeasure.ergodicDecomposition_invariantFields}
    \leanok

    Let $E$ be standard Borel and $S$ countable. If $\Theta$ is the shift group of a countable
    abelian site group admitting an increasing Følner sequence $(F_n)$ with
    $|F_n - F_n + F_n| \leq C |F_n|$ (for instance $\mathbb Z^d$ with the cubes $[0,n)^d$), and
    $\Prob_\Theta \neq \emptyset$, there is a $(\Prob_\Theta, \mathcal I)$-kernel $\pi$. For any
    countable $\Theta$ with such a kernel: $\operatorname{ex} \Prob_\Theta \neq \emptyset$; every
    $\mu \in \Prob_\Theta$ is $\mu = \int \nu\, w_\mu(d\nu)$ for a unique probability weight
    $w_\mu$ on $\operatorname{ex} \Prob_\Theta$; $\mu \mapsto w_\mu$ is an affine bijection of
    $\Prob_\Theta$ onto these weights; $w_\mu$ is the image of $\mu$ under
    $\omega \mapsto \pi(\cdot \mid \omega)$; and
    $w_\mu(\nu(A) \leq c) = \mu(\mu(A \mid \mathcal I) \leq c)$ for measurable $A$ and real $c$.
\end{theorem}
\begin{proof}
    \uses{thm:erg-14A8, thm:ext-726}
    \leanok
    The individual ergodic theorem \ref{thm:erg-14A8} gives, on a set $\Omega_0 \in \mathcal I$ of
    full measure for every $\mu \in \Prob_\Theta$, the limits of the averages of the indicators
    of a countable generating class; these limits are the distribution functions of a kernel
    $\pi(\cdot \mid \omega)$ which is a version of $\mu(\cdot \mid \mathcal I)$ for every
    $\mu \in \Prob_\Theta$, and $\pi(\cdot \mid \omega)$ is invariant and ergodic. The rest is the
    abstract extreme decomposition governed by a $(P, \mathcal A)$-kernel.
\end{proof}

\begin{corollary}[Georgii (14.11)]
    \label{cor:erg-1411}
    \uses{thm:erg-1410}
    \lean{MeasureTheory.GibbsMeasure.weight_map_of_mem_normalizer}
    \leanok

    In the setting of Theorem \ref{thm:erg-1410}, if $\tau \in T$ normalises $\Theta$ then
    $w_{\tau\mu} = \tau w_\mu$ for all $\mu \in \Prob_\Theta$.
\end{corollary}
\begin{proof}
    \leanok
    $\tau$ maps $\Prob_\Theta$ and $\operatorname{ex}\Prob_\Theta$ to themselves; apply the
    uniqueness in Theorem \ref{thm:erg-1410}.
\end{proof}

\begin{theorem}[Georgii (14.17)]
    \label{thm:erg-1417}
    \uses{thm:erg-1410, thm:erg-1415, prop:erg-149}
    \lean{MeasureTheory.GibbsMeasure.ae_mem_invariantG_of_isPAKernel_invariantFields, MeasureTheory.GibbsMeasure.weight_extremePoints_invariantG_compl, MeasureTheory.GibbsMeasure.ergodicDecomposition_invariantG}
    \leanok

    Let $\gamma$ be a specification, $\Theta$ countable with a $(\Prob_\Theta, \mathcal I)$-kernel
    $\pi$, and suppose every $\Lambda \in \mathcal S$ is moved off itself by some
    $\tau \in \Theta$. Then $\pi(\cdot \mid \omega) \in \mathcal G_\Theta(\gamma)$ for
    $\mu$-a.e.\ $\omega$, for every $\mu \in \mathcal G_\Theta(\gamma)$; $w_\mu$ is carried by
    $\operatorname{ex}\mathcal G_\Theta(\gamma)$; every $\mu \in \mathcal G_\Theta(\gamma)$ has a
    unique extreme decomposition within $\mathcal G_\Theta(\gamma)$; and $\mu \mapsto w_\mu$ is an
    affine bijection of $\mathcal G_\Theta(\gamma)$ onto the probability weights on
    $\operatorname{ex}\mathcal G_\Theta(\gamma)$.
\end{theorem}
\begin{proof}
    \leanok
    By Proposition \ref{prop:erg-149}, $\mathcal I \subseteq \mathcal T$ modulo $\mu$, so
    $\mu(\cdot \mid \mathcal I) = \mu(\mu(\cdot \mid \mathcal T) \mid \mathcal I)$ and the DLR
    equations pass to $\pi(\cdot \mid \omega)$; the rest is Theorem \ref{thm:erg-1410} restricted
    to the face $\mathcal G_\Theta(\gamma)$ of $\Prob_\Theta$ (Theorem \ref{thm:erg-1415}).
\end{proof}

\begin{corollary}[Georgii (14.18)]
    \label{cor:erg-1418}
    \uses{thm:erg-1417}
    \lean{MeasureTheory.GibbsMeasure.le_encard_extremePoints_invariantG_iff}
    \leanok

    In the setting of Theorem \ref{thm:erg-1417}, for $N \in \mathbb N$:
    $|\operatorname{ex}\mathcal G_\Theta(\gamma)| \geq N$ iff $\mathcal G_\Theta(\gamma)$
    contains $N$ linearly independent elements.
\end{corollary}
\begin{proof}
    \leanok
    Distinct extreme points are linearly independent by the uniqueness of weights; conversely
    $N$ linearly independent measures cannot lie in the span of fewer than $N$ extreme points.
\end{proof}

\begin{theorem}[Georgii (14.20)]
    \label{thm:erg-1420}
    \uses{thm:erg-1415, thm:erg-14A8, lem:erg-1419, def:sym-trans-average}
    \lean{MeasureTheory.GibbsMeasure.ae_tendsto_integral_kernel_inv_card_smul_sum_shift_of_mem_extremePoints_invariantG, MeasureTheory.GibbsMeasure.ae_tendsto_shiftAverage_weakly_of_mem_extremePoints_invariantG, Specification.ae_tendsto_shiftAverage_withLocalConvergence_of_mem_extremePoints_invariantG, MeasureTheory.GibbsMeasure.ae_tendsto_shiftAverage_withLocalConvergence_cube}
    \leanok

    Let $\Theta$ be the shift group of a countable abelian site group $S$ with an increasing
    Følner sequence $(F_n)$ satisfying $|F_n - F_n + F_n| \leq C|F_n|$, $(\Lambda_n)$ an
    increasing sequence of finite volumes, and $\mu \in \operatorname{ex}\mathcal G_\Theta(\gamma)$.
    \begin{enumerate}
    \item For every bounded measurable $f$,
      $\gamma_{\Lambda_n}\bigl(|F_n|^{-1}\sum_{i \in F_n} f \circ \theta_i \,\big|\, \omega\bigr)
      \to \mu(f)$ for $\mu$-a.e.\ $\omega$.
    \item If $E$ is compact metrizable, $S$ countable and $\gamma$ shift-invariant, then
      $|F_n|^{-1}\sum_{i \in F_n} \gamma_{\Lambda_n + i}(\cdot \mid \theta_i\omega) \to \mu$
      weakly for $\mu$-a.e.\ $\omega$.
    \item If $\gamma = \rho\lambda$ is a $\lambda$-modification of the independent specification
      of a probability measure $\lambda$ and
      $|\{i \in F_n : \delta - i \notin \Lambda_n\}|/|F_n| \to 0$ for every $\delta \in S$, the
      convergence in (2) holds in the topology of local convergence.
    \end{enumerate}
    On $\mathbb Z^d$ the cubes $F_n = \Lambda_n = [-n,n]^d$ satisfy the hypotheses.
\end{theorem}
\begin{proof}
    \leanok
    (1) is the individual ergodic theorem \ref{thm:erg-14A8} for $\mu$, trivial on $\mathcal I$,
    composed with the reversed martingale $\mu(\cdot \mid \mathcal T_{\Lambda_n})$ through Lemma
    \ref{lem:erg-1419}. (2) follows from (1) on a countable dense set of continuous functions.
    (3): for a cylinder event $A \in \mathcal F_\Delta$ the averaged kernel has density
    $\tilde\rho^n_\Delta = |F_n|^{-1}\sum_i \rho_\Delta \circ \tilde\theta_i$ relative to
    $\lambda^S$ on $\mathcal T_\Delta \otimes \mathcal F_\Delta$; by (14.22) its conditional
    expectation is $\mu$-a.s.\ equal to $\rho_\Delta$'s, and the truncated averages
    $\min(\tilde\rho^n_\Delta, M)$ converge by \ref{thm:erg-14A8} and \ref{lem:erg-1419} to a
    limit dominating $\min(\bar\rho_\Delta, M)$; letting $M \to \infty$ and using
    $\int \tilde\rho^n_\Delta = \int \bar\rho_\Delta$, a one-sided Scheffé lemma gives
    $L^1(\lambda^\Delta)$ convergence, hence (3).
\end{proof}

\begin{corollary}[Georgii (14.25)]
    \label{cor:erg-1425}
    \uses{thm:erg-1417, thm:erg-1420, def:ql-quasilocal-spec}
    \lean{MeasureTheory.GibbsMeasure.setOf_mem_invariantG_eq_closure_convexCombos}
    \leanok

    Let $\gamma$ be quasilocal and $\Theta$, $\pi$ as in Theorem \ref{thm:erg-1417}. If
    $L \subseteq \mathcal G_\Theta(\gamma)$ contains $\operatorname{ex}\mathcal G_\Theta(\gamma)$
    (for instance the set of limits of Theorem \ref{thm:erg-1420}), then
    $\mathcal G_\Theta(\gamma)$ is the closed convex hull of $L$ in the topology of local
    convergence.
\end{corollary}
\begin{proof}
    \leanok
    $\mathcal G_\Theta(\gamma)$ is closed and convex, and by Theorem \ref{thm:erg-1417} each of
    its elements is a barycentre of extreme points, approximated by finite convex combinations
    as in the density step of Theorem (7.30).
\end{proof}

\begin{proposition}[Georgii (14.21)--(14.24); Comment (14.13)]
    \label{prop:erg-1424}
    \uses{thm:erg-1420, thm:erg-14A8, thm:erg-1412}
    \lean{Specification.shiftAvgCondDensity, Specification.toReal_shiftAvgCondDensity_ae_eq_condExp, MeasureTheory.GibbsMeasure.ae_tendsto_inv_card_smul_sum_shift_of_mem_trivialOn_of_integrable, Specification.ae_tendsto_toReal_shiftAvgDensity, MeasureTheory.GibbsMeasure.continuous_toMeasure_withLocalConvergence, MeasureTheory.GibbsMeasure.invariantFields_subset_closure_extremePoints_probabilityMeasure_int}
    \leanok

    For $\gamma = \rho\lambda$ and $\mu \in \operatorname{ex}\mathcal G_\Theta(\gamma)$, the
    averaged conditional densities $\tilde\rho^n_\Delta = |\Lambda'_n|^{-1}\sum_{i \in \Lambda'_n}\rho_\Delta\circ\tilde\theta_i$
    of (14.21) have $\mu$-conditional expectation $\bar\rho_\Delta$ on
    $\mathcal T_\Delta \otimes \mathcal F_\Delta$ (14.22), and
    $\tilde\rho^n_\Delta(\omega,\zeta) \to \bar\rho_\Delta(\zeta)$ for $\mu$-a.e.\ $\omega$ and
    $\lambda^\Delta$-a.e.\ $\zeta$ along any Følner sequence, by the individual ergodic theorem for
    integrable functions ((14.23)--(14.24)). For compact $E$, local convergence is continuous into
    the weak topology and $\operatorname{ex}\Prob_\Theta$ is weakly dense in $\Prob_\Theta$ (the
    first half of (14.13)).
\end{proposition}
